\subsection{A cubic Fourier-module obstruction}
\begin{proposition}\label{prop:cubic-Fourier-module}
The fusion ring $R7\_5\_5007$ has no complex fusion-category categorification.
\end{proposition}
\begin{proof}
Write $g=b_1$, so $g^2=b_2$, and write the four remaining,
self-dual basis objects as $X_3,X_4,X_5,X_6$. Each is fixed by $g$.

The pointed subcategory generated by $g$ has trivial associator class:
the action on any one of these fixed simple objects supplies a rank-one
module category. Thus $A=1\oplus g\oplus g^2$ is a separable algebra.
The category of right $A$-modules has one simple over the free orbit
$\{1,g,g^2\}$ and three simples over each $X_i$. The dual group
$\widehat{C_3}$ acts by twisting the module structure, cyclically permuting
each triple and fixing the free-orbit simple. This action commutes with
the left action of the original category. For a primitive cube root
$\zeta$, its $\zeta$-eigenspace therefore yields a four-dimensional
representation $b_i\mapsto T_i$ of the fusion ring. In a Fourier basis,
\[
 (T_i)_{ab}=n_0+n_1\zeta+n_2\zeta^2,
 \qquad n_j\in\mathbb Z_{\ge0},\quad
 n_0+n_1+n_2=N_{i,a}^{b}, \qquad a,b\in\{3,4,5,6\}.
\]
Frobenius reciprocity for the module action gives $T_{i^*}=T_i^*$.
In particular the four matrices $T_i$ for $i=3,4,5,6$ are Hermitian.
The matrix $T_g$ is diagonal, with cube roots of unity on the diagonal.
Since $gX_i=X_i$, the identity $T_gT_i=T_i$ implies that a nontrivial
diagonal entry of $T_g$ forces the corresponding row of every $T_i$
($i=3,4,5,6$) to vanish. But a sum of three nonnegative multiples of
$1,\zeta,\zeta^2$ vanishes only when their total is divisible by three.
The entries $N_{3,3}^{4}=1$, $N_{3,4}^{3}=1$,
$N_{3,5}^{3}=2$, and $N_{3,6}^{3}=2$ consequently force $T_g=I$.

Set $T=T_3$. The following identity holds in the fusion ring, as checked
by multiplying the given integer matrices:
\[
 (X_3-9)(X_3+1)(X_3^3-X_3^2-4X_3-1)=6-3g-3g^2.
\]
Thus $(T-9I)(T+I)(T^3-T^2-4T-I)=0$.
The matrix $B=(N_{3,a}^{b})_{a,b=3}^6$ is
\[
 B=\begin{pmatrix}3&1&2&2\\1&2&2&3\\2&2&1&3\\2&3&3&3\end{pmatrix}.
\]
Its entries dominate the absolute values of the entries of $T$.
The leading principal minors of $9I-B$ are $6,41,268,47$, so
$\rho(B)<9$ and $9$ is not an eigenvalue of $T$.
The characteristic polynomial of a Hermitian matrix over
$\mathbb Q(\zeta)$ has rational coefficients; here they are algebraic
integers and hence integers. Since $x^3-x^2-4x-1$ is irreducible over
$\mathbb Q$, the only possibilities are $(x+1)^4$ and
$(x+1)(x^3-x^2-4x-1)$. Reduction modulo $(1-\zeta)$ shows
$\operatorname{tr}T\equiv\operatorname{tr}B=9\equiv0\pmod3$,
which excludes the first possibility. Consequently
\[
 \operatorname{tr}T=0,\qquad \operatorname{tr}T^2=10,
 \qquad \operatorname{tr}T^3=15.
\]
A real diagonal entry with total multiplicity $N$ has the form $N-3b$
with $0\le b\le\lfloor N/2\rfloor$. Together with the first two
moments, this forces diagonal $(0,-1,1,0)$. For an off-diagonal entry
with multiplicity $1$ or $2$, its squared norm is at least one; for
multiplicity three its squared norm is at least zero. There are four
of the former entries and two of the latter. Since the diagonal
squared norm is two, the second-moment bound is saturated: the four
entries have norm one, and the other two vanish. The only nonzero
triangle is on labels $3,4,5$. The diagonal contributions and the
edge contributions to $\operatorname{tr}T^3$ cancel, leaving six times
the real part of a product of three norm-one entries. Its absolute
value is at most six, contradicting $\operatorname{tr}T^3=15$.
The certificate checker also exhausts all 81 possible matrices after
these bounds, obtaining cubic traces only $-3$ and $6$.

\end{proof}
